\documentclass[10pt,letterpaper]{article}

\usepackage[utf8]{inputenc}
\usepackage[T1]{fontenc}
\usepackage[spanish,es-nodecimaldot]{babel}
\usepackage[letterpaper,top=2.0cm,bottom=2.0cm,left=2.55cm,right=2.55cm]{geometry}
\usepackage{graphicx,float,booktabs,array,tabularx,amsmath,amssymb}
\usepackage{xcolor,hyperref,enumitem,caption,microtype}
\usepackage{ragged2e}

\setlength{\parindent}{0pt}
\setlength{\parskip}{3.5pt}
\setlist{nosep,leftmargin=*}
\renewcommand{\arraystretch}{1.10}
\captionsetup{font=small,labelfont=bf}
\hypersetup{colorlinks=true,linkcolor=black,urlcolor=blue,citecolor=black}
\newcolumntype{Y}{>{\RaggedRight\arraybackslash}X}
\newcommand{\pct}[1]{#1\%}

\begin{document}

% =========================================================
% ENCABEZADO INSTITUCIONAL (formato de main (1).tex)
% =========================================================
\begin{minipage}[t]{0.24\textwidth}
    \vspace{0pt}
    \centering
    \IfFileExists{logo javeriana.png}{\includegraphics[width=3.05cm]{logo javeriana.png}}{}
\end{minipage}
\hfill
\begin{minipage}[t]{0.72\textwidth}
    \vspace{0pt}
    \begin{flushright}
        {\large\textbf{\textit{Facultad de Ingeniería}}}\\[0.15cm]
        {\Large\textit{Métodos y aplicaciones de analítica I}}\\[0.10cm]
        {\small\textit{Taller 01 -- Segundo Semestre 2026}}
    \end{flushright}
\end{minipage}

\vspace{0.45cm}

\begin{center}
    {\large\textbf{Sistemas de recomendación: reglas de asociación en Instacart}}\\[0.30cm]
    {\normalsize Nestor Fandiño\textsuperscript{a,c}}\\[0.12cm]
    {\normalsize Diego Alejandro Cortes Suarez\textsuperscript{a,c}}\\[0.12cm]
    {\normalsize Andres Rodriguez Torres\textsuperscript{b,c}}\\[0.18cm]
    {\small
    \textsuperscript{a}\textit{Estudiantes de la Maestría en Inteligencia Artificial}\\[-0.02cm]
    \textsuperscript{b}\textit{Profesor, Departamento de Ingeniería de Sistemas}\\[-0.02cm]
    \textsuperscript{c}\textit{Pontificia Universidad Javeriana, Bogotá, Colombia}}
\end{center}

\vspace{0.30cm}

\section*{1. ¿Qué hicimos?}

Tomamos 50\,000 órdenes históricas de Instacart y tratamos cada orden como una canasta. Después buscamos patrones del tipo $n\rightarrow1$: uno, dos o tres productos/pasillos que aparecen en la canasta y un único elemento que se puede recomendar.

Ejemplo: si una canasta tiene \texttt{dry pasta} y \texttt{packaged vegetables fruits}, podemos evaluar si conviene sugerir \texttt{pasta sauce}. No estamos diciendo que la pasta o las verduras \emph{causen} la compra de salsa; solo que, en los datos, aparecen juntas más de lo esperado.

Trabajamos dos niveles:
\begin{itemize}
    \item \textbf{Producto $\rightarrow$ producto:} más específico, pero con menos casos por combinación.
    \item \textbf{Pasillo $\rightarrow$ pasillo:} más general y más útil para decisiones de tienda, porque agrupa artículos similares.
\end{itemize}

\section*{2. La métrica que guía las decisiones: lift}

Para una regla $A\Rightarrow B$, el \textit{lift} compara la probabilidad de comprar $B$ cuando ya está $A$ en la canasta con la probabilidad normal de comprar $B$:
\[
\textit{lift}(A\Rightarrow B)=\frac{P(B\mid A)}{P(B)}.
\]

\begin{center}
\begin{tabular}{cl}
\toprule
Lift & Lectura práctica \\
\midrule
$>1$ & Hay asociación positiva: $B$ aparece más de lo normal cuando está $A$. \\
$=1$ & No hay ganancia clara: la compra de $A$ no cambia la probabilidad de $B$. \\
$<1$ & La combinación aparece menos de lo esperado. \\
\bottomrule
\end{tabular}
\end{center}

El soporte y la confianza siguen siendo controles: el soporte dice cuántas veces vemos la regla y la confianza qué proporción de las canastas con $A$ también tiene $B$. Pero el \textit{lift} es el centro del análisis: evita recomendar un producto solo porque es popular. Si el \textit{lift} es cercano a 1, la recomendación no aporta mucho.

\section*{3. ¿Se pueden usar reglas $n\rightarrow1$?}

Sí. Las canastas tienen tamaño suficiente para explorar antecedentes de hasta tres elementos:
\begin{center}
\begin{tabular}{lrrr}
\toprule
Nivel & $1\rightarrow1$ & $2\rightarrow1$ & $3\rightarrow1$ \\
\midrule
Productos & \pct{95,24} & \pct{89,36} & \pct{82,98} \\
Pasillos & \pct{93,87} & \pct{86,01} & \pct{77,17} \\
\bottomrule
\end{tabular}
\end{center}

Usamos $1\rightarrow1$ y $2\rightarrow1$ como resultados principales porque son fáciles de activar en un carrito y de explicar. Las reglas $3\rightarrow1$ también funcionan, pero requieren demasiado contexto para una primera recomendación.

\section*{4. Resultados: reglas que sí sugieren una acción}

Las siguientes reglas se descubrieron en órdenes históricas y se revisaron de nuevo en \texttt{train}, un conjunto diferente. El \textit{lift} de la tabla es el de validación: la señal se mantiene fuera de las órdenes que se usaron para descubrirla.

\begin{table}[H]
\centering
\caption{Reglas de pasillo recomendadas.}
\scriptsize
\begin{tabularx}{\linewidth}{YYrrrY}
\toprule
Si la canasta contiene... & Sugerir... & Soporte & Confianza & Lift & ¿Qué hacer? \\
\midrule
\texttt{pasta sauce} & \texttt{dry pasta} & \pct{2,216} & \pct{32,97} & \textbf{4,54} & Sugerir pasta seca en el carrito; crear una exhibición conjunta. \\
\texttt{dry pasta + packaged vegetables fruits} & \texttt{pasta sauce} & \pct{1,268} & \pct{32,47} & \textbf{4,83} & Mostrar salsa como recomendación contextual: es la regla más fuerte de la tabla. \\
\texttt{fresh vegetables + pasta sauce} & \texttt{dry pasta} & \pct{1,495} & \pct{33,13} & \textbf{4,56} & Activar sugerencia de pasta cuando la canasta ya parece una comida completa. \\
\texttt{frozen pizza} & \texttt{frozen meals} & \pct{1,229} & \pct{25,08} & \textbf{3,05} & Probar una recomendación complementaria en congelados. \\
\texttt{canned meals beans + fresh fruits + fresh vegetables} & \texttt{canned jarred vegetables} & \pct{1,265} & \pct{29,99} & \textbf{3,90} & Usarla como regla avanzada; no como primera recomendación. \\
\bottomrule
\end{tabularx}
\end{table}

\textbf{Cómo leer el mejor resultado.} La regla \texttt{dry pasta + packaged vegetables fruits} $\Rightarrow$ \texttt{pasta sauce} tiene un \textit{lift} de 4,83. Esto significa que, cuando se observan pasta seca y verduras empacadas en la misma canasta, la compra de salsa es aproximadamente 4,83 veces más probable que para una canasta cualquiera. Por eso es una buena candidata para mostrar salsa, siempre que el producto no esté ya en el carrito.

\section*{5. Producto versus pasillo}

También encontramos reglas muy fuertes por producto. Por ejemplo, \texttt{Banana + Limes} $\Rightarrow$ \texttt{Large Lemon} tuvo \textit{lift} de 5,14. Estas reglas pueden servir para recomendaciones muy específicas, pero suelen tener menos observaciones: en ese caso fueron 131 órdenes de las 50\,000 analizadas.

Por esta razón, proponemos usar los pasillos para la recomendación base y dejar las reglas de producto como una segunda capa, cuando el carrito ya contiene los artículos exactos del antecedente. Así evitamos depender de relaciones de nicho, aunque tengan un \textit{lift} alto.

\section*{6. Qué haríamos con estos resultados}

\begin{enumerate}
    \item \textbf{Tienda virtual:} cuando se cumpla un antecedente, mostrar un solo producto del pasillo consecuente. Priorizar reglas con mayor \textit{lift}; si hay empate, usar mayor soporte. No mostrar productos que ya estén en el carrito.
    \item \textbf{Tienda física:} usar las reglas de pasta, salsa y verduras para señalización, promociones o exhibiciones cercanas. El resultado sirve como hipótesis para una prueba, no como orden automática de mover productos.
    \item \textbf{Medición:} comparar una versión con recomendación contra otra sin recomendación y medir adiciones al carrito, compras y margen. Un \textit{lift} alto en datos históricos es una buena señal, pero no prueba por sí solo que la recomendación aumente ventas.
\end{enumerate}

\section*{7. Cierre}

El resultado útil no es una lista larga de combinaciones, sino una regla de decisión: recomendamos cuando el \textit{lift} es claramente mayor que 1, la regla aparece suficientes veces y tiene sentido comercial. En esta muestra, las reglas por pasillo $1\rightarrow1$ y $2\rightarrow1$ son las más adecuadas para empezar: tienen \textit{lift} entre 3,05 y 4,83, son fáciles de entender y se mantuvieron al validarlas en \texttt{train}.

\vspace{0.25cm}
\begin{thebibliography}{9}
\bibitem{instacart} Instacart (2017). \textit{The Instacart Online Grocery Shopping Dataset 2017}. \url{https://www.instacart.com/datasets/grocery-shopping-2017}.
\bibitem{agrawal} R. Agrawal y R. Srikant (1994). Fast algorithms for mining association rules. \textit{Proceedings of the 20th International Conference on Very Large Data Bases}.
\end{thebibliography}

\end{document}
